%% toc: P15 160 160.5,160.6,160.7,160.8
\chapter{Essential Bounds, the Moment Method and Urn Parameters}
\label{sec:p15-w02c-p15-012-bounds-and-urns}

Two results in the preceding chapters depend on parameter conditions that
deserve their own statements. The first is the consequence that Durrett
draws from the essential bound of a random variable: once the bound is
finite, the moment-growth hypothesis of the method of moments holds
automatically \cite[p.~143, Exercise~3.3.24 and Theorem~3.3.26]{DurrettPTEAuthor2019}.
The second is Feller's identification of the P\'olya urn's mixing law
through its all-success probabilities, where the beta law exists only when
both initial weights are positive
\cite[vol.~II, VII.4, pp.~229--230]{Feller1971ProbabilityII}.
This chapter states both with every hypothesis explicit, proves them in
full and records what happens when a hypothesis is removed.

\paragraph{Inputs and conclusions.}
The first section takes a finite real random variable \(X\) and its
essential bound \(\Lambda\) as in
\cref{thm:p15-criteria-compact-moments}, and concludes that
\(\Lambda<\infty\) gives the growth hypothesis of
\cref{thm:p15-criteria-moment-method} for the moments of \(X\), so that
moment convergence to those moments is weak convergence to the law of
\(X\). The second section takes urn weights \(b,r\geq0\), \(b+r>0\), and
reinforcement \(c>0\), and concludes which probability on \([0,1]\) is the
mixing law of the black-draw indicators: \(\operatorname{beta}(b/c,r/c)\)
when \(b,r>0\), and a point mass when one weight is zero.

\section{The essential bound and the moment-growth hypothesis}
\label{sec:p15-w02c-p15-012-essential-bound}

\begin{hbconvention}{def:p15-w02c-p15-012-moment-notation}{Moments, absolute moments and the essential bound}
For a finite real random variable \(X\), the notation of
\cref{def:p15-criteria-moment-problem} and
\cref{thm:p15-criteria-compact-moments} is retained:
\[
 \nu_k=\mathbb E|X|^k\in[0,\infty],\qquad
 \Lambda=\sup\{x:\mathbb P(|X|<x)<1\}\in[0,\infty].
\]
When every \(\nu_k\) is finite, \(m_k=\mathbb EX^k\) denotes the
\(k\)-th moment; this is the moment sequence that appears in the
hypothesis of \cref{thm:p15-criteria-moment-method}.
Durrett writes \(G(x)=\mathbb P(|X|<x)\) and \(\Lambda=\sup\{x:G(x)<1\}\)
for the same bound, and writes \(\mu_k\) for the limit moments in
his Theorem~3.3.26 \cite[p.~143]{DurrettPTEAuthor2019}; here they are
\(m_k\). Since \(2k\) is even, \(m_{2k}=\nu_{2k}\).
\end{hbconvention}

\begin{hbcorollary}{thm:p15-w02c-p15-012-bound-gives-growth}{proved=annex:p15-w02c-p15-012-bound-gives-growth}{A finite essential bound gives the moment-method hypothesis}
Let \(X\) be a finite real random variable with \(\Lambda<\infty\).
Then:
\begin{enumerate}
\item \(|X|\leq\Lambda\) almost surely, every \(\nu_k\) is finite, and
      \(\nu_k\leq\Lambda^k\) for every \(k\geq1\);
\item \(m_{2k}^{1/(2k)}\to\Lambda\) and
      \[
       0\leq\frac{m_{2k}^{1/(2k)}}{2k}\leq\frac{\Lambda}{2k}
       \quad(k\geq1),\qquad\text{so}\qquad
       \limsup_{k\to\infty}\frac{m_{2k}^{1/(2k)}}{2k}=0<\infty;
      \]
\item the law of \(X\) is the only probability law on \(\mathbb R\) with
      the moments \(m_k\), and the hypothesis of
      \cref{thm:p15-criteria-moment-method} holds for the sequence
      \(m_k\); consequently, if probability laws \(\mu_n\) have all finite
      absolute moments and \(\int x^k\,\mu_n(dx)\to m_k\) for every positive
      integer \(k\), then \(\mu_n\) converges weakly to the law of \(X\).
\end{enumerate}
This is the clause of Durrett's Exercise~3.3.24 that follows the limit
\(\nu_k^{1/k}\to\Lambda\): the assumption of his Theorem~3.3.26 holds
when \(\Lambda<\infty\) \cite[p.~143]{DurrettPTEAuthor2019}. The limit
itself is the first assertion of \cref{thm:p15-criteria-compact-moments}.
\index{moment problem!essential bound}\index{method of moments}
\end{hbcorollary}

The bound in item~2 is stronger than the hypothesis requires: the
normalized root tends to zero, so the analytic determinacy theorem
\cref{thm:p15-criteria-analytic-determinacy} applies with \(r=0\).
Determinacy of a compactly supported
law was proved independently in \cref{thm:p15-criteria-compact-moments};
item~3 obtains it a second way and adds the convergence conclusion.

\begin{hbexample}{def:p15-w02c-p15-012-uniform-moments}{Uniform grids converge to the uniform law by moments}
Let \(X\) be uniform on \([0,1]\), so \(\Lambda=1\) and
\(m_k=1/(k+1)\). Then \(m_{2k}^{1/(2k)}=(2k+1)^{-1/(2k)}\to1=\Lambda\),
and \(m_{2k}^{1/(2k)}/(2k)\leq1/(2k)\).
Let \(\mu_n\) put mass \(1/n\) at each of \(j/n\), \(j=1,\ldots,n\).
Each \(\mu_n\) has finite support, hence all finite absolute moments, and
\[
 \frac1{k+1}\leq\int x^k\,\mu_n(dx)=\frac1n\sum_{j=1}^n\Big(\frac jn\Big)^k
 \leq\frac1{k+1}+\frac1n .
\]
So the moments converge to \(m_k\) for every \(k\), and
\cref{thm:p15-w02c-p15-012-bound-gives-growth} shows
that \(\mu_n\) converges weakly to the uniform law on \([0,1]\), with no direct computation of
distribution functions. The annex proves the displayed bounds as part of
\cref{annex:p15-w02c-p15-012-bound-gives-growth}.
\end{hbexample}

\begin{hbwarning}{def:p15-w02c-p15-012-infinite-bound}{An infinite essential bound decides nothing}
The implication runs in one direction only. When \(\Lambda=\infty\), the
limit \(\nu_k^{1/k}\to\infty\) of \cref{thm:p15-criteria-compact-moments}
holds whether or not the growth hypothesis holds.
By \cref{thm:p15-criteria-generalized-gamma} with \(\alpha=0\), if \(|X|\)
has density proportional to \(e^{-x^\lambda}\) on \((0,\infty)\), then
\(\Lambda=\infty\) for every \(\lambda>0\), while the growth hypothesis
holds for \(\lambda=2\) (for instance, \(X\) centred normal with variance \(1/2\)) and
fails for \(\lambda=1/2\). A finite essential bound is therefore a
sufficient condition for the hypothesis of
\cref{thm:p15-criteria-moment-method}, not a necessary one; and failure of
that hypothesis does not by itself show indeterminacy, as
\cref{thm:p15-criteria-generalized-gamma} records.
\end{hbwarning}

\section{P\'olya's urn in Feller's moment route: every parameter explicit}
\label{sec:p15-w02c-p15-012-urn-parameters}

Feller's example starts from an urn with \(b\) black and \(r\) red balls;
after each draw the ball is returned together with \(c\) balls of the
colour drawn, and the probability \(c_n\) of black on each of the first
\(n\) draws is computed as a product
\cite[vol.~II, VII.4, Example~(a), pp.~229--230]{Feller1971ProbabilityII}.
Feller then identifies the distribution with moments \(c_n\) as the beta
distribution with parameters \(b/c\) and \(r/c\). That identification
uses both \(b>0\) and \(r>0\): the beta normalizing integral is finite
exactly when both parameters are positive. The statements below therefore
separate the case of two positive weights from the cases with one weight
zero. The argument uses only the positivity of the weights, so real
weights are allowed; integer weights are Feller's ball counts.

\begin{hbdefinition}{def:p15-w02c-p15-012-urn}{The urn with explicit weights}
Fix real \(b\geq0\), \(r\geq0\) with \(b+r>0\), and real \(c>0\).
Let \(U_1,U_2,\ldots\) be independent and uniform on \((0,1)\).
Put \(S_0=0\) and, recursively for \(n\geq0\),
\[
 p_n=\frac{b+cS_n}{b+r+cn},\qquad
 X_{n+1}=\mathbf 1_{\{U_{n+1}\leq p_n\}},\qquad
 S_{n+1}=S_n+X_{n+1}.
\]
Thus \(X_j=1\) records a black draw, \(S_n\) counts the black draws among
the first \(n\), and \(p_n\) is the black weight divided by the total
weight after \(n\) draws. The denominator is positive because \(b+r>0\)
and \(c>0\). Put \(c_0=1\) and
\[
 c_n=\mathbb P(X_1=\cdots=X_n=1)\quad(n\geq1).
\]
For \(a\geq0\) the rising product is \((a)_0=1\),
\((a)_n=a(a+1)\cdots(a+n-1)\). For \(\alpha,\beta>0\),
\[
 B(\alpha,\beta)=\int_0^1\theta^{\alpha-1}(1-\theta)^{\beta-1}\,d\theta,
\]
and \(\operatorname{beta}(\alpha,\beta)\) is the probability on \([0,1]\)
with density \(\theta^{\alpha-1}(1-\theta)^{\beta-1}/B(\alpha,\beta)\) on
\((0,1)\).
\index{P\'olya urn!parameter boundaries}\index{beta distribution}
\end{hbdefinition}

\begin{hblemma}{thm:p15-w02c-p15-012-urn-words}{proved=annex:p15-w02c-p15-012-urn-words}{Word probabilities and exchangeability for all admissible weights}
In the urn of \cref{def:p15-w02c-p15-012-urn}, every binary word
\(x_1,\ldots,x_n\) with \(k=\sum_jx_j\) black entries satisfies
\begin{align*}
 \mathbb P(X_1=x_1,\ldots,X_n=x_n)
   &=\frac{\prod_{i=0}^{k-1}(b+ic)\,\prod_{i=0}^{n-k-1}(r+ic)}
         {\prod_{i=0}^{n-1}(b+r+ic)}\\
   &=\frac{(b/c)_k\,(r/c)_{n-k}}{((b+r)/c)_n}.
\end{align*}
The value depends on the word only through \(k\); hence for every \(n\)
and every permutation \(\pi\) of \(\{1,\ldots,n\}\), the vectors
\((X_{\pi1},\ldots,X_{\pi n})\) and \((X_1,\ldots,X_n)\) have the same law,
and \(X_1,X_2,\ldots\) is an infinite exchangeable binary sequence.
In particular
\[
 c_n=\frac{b(b+c)\cdots(b+(n-1)c)}{(b+r)(b+r+c)\cdots(b+r+(n-1)c)}
    =\frac{(b/c)_n}{((b+r)/c)_n}.
\]
\end{hblemma}

\begin{hbtheorem}{thm:p15-w02c-p15-012-urn-beta}{proved=annex:p15-w02c-p15-012-urn-beta}{Feller's urn identification with both weights positive}
Assume \(b>0\), \(r>0\) and \(c>0\), and put \(\alpha=b/c\), \(\beta=r/c\).
Then:
\begin{enumerate}
\item \(B(\alpha,\beta)\) is finite and positive, and the moments of
      \(\operatorname{beta}(\alpha,\beta)\) are
      \begin{align*}
       \int\theta^n\,\operatorname{beta}(\alpha,\beta)(d\theta)
         &=\frac{B(\alpha+n,\beta)}{B(\alpha,\beta)}
          =\frac{(\alpha)_n}{(\alpha+\beta)_n}\\
         &=\frac{\Gamma(\alpha+n)\Gamma(\alpha+\beta)}
               {\Gamma(\alpha+\beta+n)\Gamma(\alpha)}
          =c_n ;
      \end{align*}
\item the mixing probability \(\lambda\) on \([0,1]\) that
      \cref{thm:p15-feller-binary-mixture} attaches to the exchangeable
      sequence \(X_1,X_2,\ldots\) is \(\operatorname{beta}(\alpha,\beta)\);
\item every word of length \(n\) with \(k\) black entries has probability
      \(B(\alpha+k,\beta+n-k)/B(\alpha,\beta)\), and the number of black
      draws among the first \(n\) has
      \[
       \mathbb P(S_n=k)=\binom nk\frac{B(\alpha+k,\beta+n-k)}{B(\alpha,\beta)}
                     \quad(0\leq k\leq n).
      \]
\end{enumerate}
\index{P\'olya urn!beta mixing law}
\end{hbtheorem}

\begin{hbtheorem}{thm:p15-w02c-p15-012-urn-one-colour}{proved=annex:p15-w02c-p15-012-urn-one-colour}{One initial weight zero: the mixing law is a point mass}
Let \(c>0\).
\begin{enumerate}
\item If \(b>0=r\), then \(X_j=1\) for every \(j\) surely, \(c_n=1\) for
      every \(n\geq0\), and the mixing probability of
      \cref{thm:p15-feller-binary-mixture} is \(\delta_1\). The product and
      rising-product formulas for \(c_n\) remain valid and give \(1\), but
      \(B(b/c,0)\) diverges, so no beta law with parameters \(b/c\) and
      \(0\) exists.
\item If \(b=0<r\), then \(X_j=0\) for every \(j\) surely, \(c_n=0\) for
      every \(n\geq1\), and the mixing probability is \(\delta_0\). The
      product formula gives \(c_n=0\) through its first factor \(b=0\);
      \(B(0,r/c)\) diverges, and the gamma expression is undefined because
      \(\Gamma\) has a pole at \(0\).
\end{enumerate}
In both cases the identification by Hausdorff uniqueness still holds; only
the beta parametrization fails.
\end{hbtheorem}

\begin{hbexample}{def:p15-w02c-p15-012-urn-example}{Two black, one red, one added ball, and the one-colour contrast}
Take \(b=2\), \(r=1\), \(c=1\), so \(\alpha=2\), \(\beta=1\).
\Cref{thm:p15-w02c-p15-012-urn-words} shows
\(c_n=(2)_n/(3)_n=2/(n+2)\). The law \(\operatorname{beta}(2,1)\) has
density \(2\theta\) on \((0,1)\), with moments
\(\int_0^1 2\theta^{n+1}\,d\theta=2/(n+2)\), as
\cref{thm:p15-w02c-p15-012-urn-beta} asserts.
The ordered word black, black, red has probability
\((2\cdot3\cdot1)/(3\cdot4\cdot5)=1/10\), and also
\(B(4,2)/B(2,1)=(1/20)/(1/2)=1/10\).
This is the word of \cref{thm:p15-polya-worked-prediction} with green
renamed black.
Now remove the red ball: \(b=2\), \(r=0\), \(c=1\). Every draw is black,
\(c_n=(2)_n/(2)_n=1\), and the mixing law is \(\delta_1\); the candidate
density \(\theta(1-\theta)^{-1}\) is not integrable near \(1\).
\end{hbexample}

\begin{hbwarning}{def:p15-w02c-p15-012-urn-limitations}{What the parameter conditions exclude}
The beta identification needs \(c>0\) as well as \(b,r>0\): with \(c=0\)
the parameters \(b/c,r/c\) are undefined. That process draws independent
Bernoulli variables with success probability \(b/(b+r)\), so its mixing
law is the point mass at \(b/(b+r)\); this case and the limits
\(c\downarrow0\) and \(c\to\infty\) are proved in
\cref{thm:p15-polya-parameter-boundaries}. With \(b=r=0\) the first draw
is undefined, so \cref{def:p15-w02c-p15-012-urn} requires \(b+r>0\).
The theorems here concern the all-success moments and the mixing law;
the almost-sure limit of the black fraction, with the same laws, is
\cref{thm:p15-polya-beta-limit} for positive weights and
\cref{thm:p15-polya-parameter-boundaries} for one zero weight.
\end{hbwarning}

\section{Relationships}
\label{sec:p15-w02c-p15-012-relationships}

\Cref{thm:p15-w02c-p15-012-bound-gives-growth} completes Durrett's
Exercise~3.3.24 together with \cref{thm:p15-criteria-compact-moments},
which proves the limit \(\nu_k^{1/k}\to\Lambda\) and compact determinacy.
It connects that limit to the method of moments,
\cref{thm:p15-criteria-moment-method}, and to analytic determinacy,
\cref{thm:p15-criteria-analytic-determinacy}. With it, the list of
Section~3.3 exercises in \cref{def:p15-criteria-completion-boundary} and
\cref{def:p15-characteristic-next-boundary} has Exercise~3.3.24 in full.

\Cref{thm:p15-w02c-p15-012-urn-beta} is
\cref{thm:p15-feller-polya-moments} with its parameter conditions written
out: that theorem is to be read with black weight \(b>0\), red weight
\(r>0\) and \(c>0\). \Cref{thm:p15-w02c-p15-012-urn-one-colour} supplies
the cases with one zero weight inside the same moment route, using
\cref{thm:p15-feller-binary-mixture} and Hausdorff uniqueness,
\cref{thm:p15-feller-hausdorff}. The direct urn route, with the green
colour in place of black, is \cref{sec:p15-polya-urns}.

\paragraph{Sources.}
Durrett, \emph{Probability: Theory and Examples}, version~5 author PDF,
Exercise~3.3.24, Theorem~3.3.25 and Theorem~3.3.26, pp.~142--143
\cite[pp.~142--143]{DurrettPTEAuthor2019}.
Feller, \emph{An Introduction to Probability Theory and Its Applications},
vol.~II, second edition, Section~VII.4, Example~(a) and display~(4.7),
pp.~229--230 \cite[vol.~II, VII.4, pp.~229--230]{Feller1971ProbabilityII}.
